\begin{lexlean}{PointwiseGraph}
\useglossary{lexlean.std.bool@1.2.0}
\useglossary{lexlean.std.nat@1.2.0}
\importmodule{TensorIndex}
\importmodule{MatrixGraph}
\importmodule{TensorTypes}
\title{Natural number}
\begin{semanticmodule}
structure{
  name: TensorDescriptor,
  type_parameters: [],
  parameters: [],
  fields: [
    {name: ty, type: Nat},
    {name: n0, type: Nat},
    {name: n1, type: Nat},
    {name: n2, type: Nat},
    {name: n3, type: Nat},
    {name: s0, type: Nat},
    {name: s1, type: Nat},
    {name: s2, type: Nat},
    {name: s3, type: Nat},
    {name: span, type: Nat}
  ]
}

structure{
  name: IndexContext,
  type_parameters: [],
  parameters: [],
  fields: [
    {name: u0, type: Nat},
    {name: n1, type: Nat},
    {name: n2, type: Nat},
    {name: n3, type: Nat},
    {name: s0, type: Nat},
    {name: s1, type: Nat},
    {name: s2, type: Nat},
    {name: s3, type: Nat}
  ]
}

inductive{
  name: CopyMode,
  type_parameters: [],
  parameters: [],
  constructors: [
    {name: linear, fields: []},
    {name: repeated, fields: []},
    {name: concatenated, fields: [Nat, IndexContext]}
  ]
}

inductive{
  name: PointwiseDecision,
  type_parameters: [],
  parameters: [],
  constructors: [
    {name: unsupported, fields: []},
    {name: invalidShape, fields: []},
    {name: execute, fields: [Nat, Bool, Nat, Nat, Nat, Nat, Nat]},
    {name: product, fields: [Nat, Nat, Nat, Nat, Nat, Nat, Nat, Nat, Nat, Nat]},
    {name: copyBlocks, fields: [Nat, Nat, IndexContext, IndexContext, CopyMode]}
  ]
}

definition{
  name: scalarKind,
  parameters: [operation : Nat, variant : Nat],
  result: (Option Nat),
  body: (if (operation =? 2) then (ctor Option.some [Nat] 0) else (if (operation =? 5) then (ctor Option.some [Nat] 1) else (if (operation =? 6) then (ctor Option.some [Nat] 2) else (if (operation =? 7) then (ctor Option.some [Nat] 3) else (if (operation =? 9) then (ctor Option.some [Nat] 4) else (if (operation =? 8) then (ctor Option.some [Nat] 5) else (if (operation =? 89) then (if (variant =? 2) then (ctor Option.some [Nat] 6) else (if (variant =? 0) then (ctor Option.some [Nat] 7) else (ctor Option.none [Nat]))) else (ctor Option.none [Nat])))))))),
  executable: true
}

definition{
  name: empty,
  parameters: [n0 : Nat, n1 : Nat, n2 : Nat, n3 : Nat],
  result: Bool,
  body: ((((n0 =? 0) || (n1 =? 0)) || (n2 =? 0)) || (n3 =? 0)),
  executable: true
}

definition{
  name: repeatValid,
  parameters: [n0 : Nat, n1 : Nat, n2 : Nat, n3 : Nat, r0 : Nat, r1 : Nat, r2 : Nat, r3 : Nat],
  result: Bool,
  body: (if (empty r0 r1 r2 r3) then (empty n0 n1 n2 n3) else (((((prim remainder n0 r0 1 : Nat) =? 0) && ((prim remainder n1 r1 1 : Nat) =? 0)) && ((prim remainder n2 r2 1 : Nat) =? 0)) && ((prim remainder n3 r3 1 : Nat) =? 0))),
  executable: true
}

definition{
  name: planPointwise,
  parameters: [
    operation : Nat,
    variant : Nat,
    leftType : Nat,
    rightType : Nat,
    outputType : Nat,
    rightPresent : Bool,
    n0 : Nat,
    n1 : Nat,
    n2 : Nat,
    n3 : Nat,
    r0 : Nat,
    r1 : Nat,
    r2 : Nat,
    r3 : Nat,
    o0 : Nat,
    o1 : Nat,
    o2 : Nat,
    o3 : Nat
  ],
  result: PointwiseDecision,
  body: (match (scalarKind operation variant)
    | Option.none => (if rightPresent then (match (MatrixGraph::planMatrix operation leftType rightType outputType n0 n1 n2 n3 r0 r1 r2 r3 o0 o1 o2 o3)
      | MatrixGraph::MatrixDecision.unsupported => (ctor PointwiseDecision.unsupported)
      | MatrixGraph::MatrixDecision.invalidShape => (ctor PointwiseDecision.invalidShape)
      | MatrixGraph::MatrixDecision.product rows inner columns b2 b3 ratio2 ratio3 leftCount rightCount outputCount => (ctor PointwiseDecision.product rows inner columns b2 b3 ratio2 ratio3 leftCount rightCount outputCount)) else (ctor PointwiseDecision.invalidShape))
    | Option.some kind => (let unary : Bool := (4 <=? kind); (if (((((leftType =? 0) && (outputType =? 0)) && ((n0 =? o0) && (n1 =? o1))) && ((n2 =? o2) && (n3 =? o3))) && (unary || ((rightPresent && (rightType =? 0)) && (repeatValid n0 n1 n2 n3 r0 r1 r2 r3)))) then (ctor PointwiseDecision.execute kind unary n0 n1 n2 n3 (prim multiply (prim multiply n0 n1 : Nat) (prim multiply n2 n3 : Nat) : Nat)) else (ctor PointwiseDecision.invalidShape)))),
  executable: true
}

definition{
  name: rightOffset,
  parameters: [
    i0 : Nat,
    i1 : Nat,
    i2 : Nat,
    i3 : Nat,
    r0 : Nat,
    r1 : Nat,
    r2 : Nat,
    r3 : Nat,
    s0 : Nat,
    s1 : Nat,
    s2 : Nat,
    s3 : Nat
  ],
  result: Nat,
  body: (TensorIndex::elementOffset (prim remainder i0 r0 0 : Nat) (prim remainder i1 r1 0 : Nat) (prim remainder i2 r2 0 : Nat) (prim remainder i3 r3 0 : Nat) s0 s1 s2 s3),
  executable: true,
  production: {effects: [overflow], targets: [rust-core]}
}

theorem{
  name: scalarBroadcastAdmitted,
  parameters: [],
  statement: ((planPointwise 2 0 0 0 0 true 2 3 2 2 1 1 1 1 2 3 2 2) = (ctor PointwiseDecision.execute 0 false 2 3 2 2 24)),
  proof: reflexivity{}
}

theorem{
  name: missingBinaryOperandRefused,
  parameters: [],
  statement: ((planPointwise 2 0 0 0 0 false 1 1 1 1 1 1 1 1 1 1 1 1) = (ctor PointwiseDecision.invalidShape)),
  proof: reflexivity{}
}

theorem{
  name: scalarProductAdmitted,
  parameters: [],
  statement: ((planPointwise 28 0 0 0 0 true 1 1 1 1 1 1 1 1 1 1 1 1) = (ctor PointwiseDecision.product 1 1 1 1 1 1 1 1 1 1)),
  proof: reflexivity{}
}

inductive{
  name: GraphRequest,
  type_parameters: [],
  parameters: [],
  constructors: [
    {name: plan, fields: [Nat, Nat, TensorDescriptor, (Option TensorDescriptor), TensorDescriptor]},
    {name: offsets, fields: [Nat, IndexContext, IndexContext, CopyMode]}
  ]
}

inductive{
  name: GraphResponse,
  type_parameters: [],
  parameters: [],
  constructors: [
    {name: planned, fields: [PointwiseDecision]},
    {name: offsets, fields: [Bool, Nat, Nat]}
  ]
}

definition{
  name: copyOperation,
  parameters: [operation : Nat],
  result: Bool,
  body: (((operation =? 1) || (operation =? 33)) || (operation =? 34)),
  executable: true
}

definition{
  name: blockCount,
  parameters: [tensor : TensorDescriptor, block : Nat],
  result: Nat,
  body: (prim multiply (prim multiply (prim quotient (project tensor n0) block 0 : Nat) (project tensor n1) : Nat) (prim multiply (project tensor n2) (project tensor n3) : Nat) : Nat),
  executable: true
}

definition{
  name: indexContext,
  parameters: [tensor : TensorDescriptor, block : Nat],
  result: IndexContext,
  body: (record IndexContext {u0 := (prim quotient (project tensor n0) block 0 : Nat), n1 := (project tensor n1), n2 := (project tensor n2), n3 := (project tensor n3), s0 := (project tensor s0), s1 := (project tensor s1), s2 := (project tensor s2), s3 := (project tensor s3)}),
  executable: true
}

definition{
  name: planStorage,
  parameters: [
    operation : Nat,
    left : TensorDescriptor,
    rightPresent : Bool,
    output : TensorDescriptor
  ],
  result: PointwiseDecision,
  body: (if (((project left ty) =? (project output ty)) && (if (operation =? 33) then rightPresent else true)) then (match (TensorTypes::tensorTypeStorage (project left ty))
    | TensorTypes::TypeStorageDecision.unknownType => (ctor PointwiseDecision.unsupported)
    | TensorTypes::TypeStorageDecision.knownType spec => (let block : Nat := (project spec blockSize); (let count : Nat := (blockCount left block); (if ((((prim remainder (project left n0) block 1 : Nat) =? 0) && ((prim remainder (project output n0) block 1 : Nat) =? 0)) && (count =? (blockCount output block))) then (ctor PointwiseDecision.copyBlocks count (project spec typeSize) (indexContext left block) (indexContext output block) (ctor CopyMode.linear)) else (ctor PointwiseDecision.invalidShape))))) else (ctor PointwiseDecision.invalidShape)),
  executable: true
}

definition{
  name: planRepeat,
  parameters: [left : TensorDescriptor, output : TensorDescriptor],
  result: PointwiseDecision,
  body: (if (((((project left ty) =? 0) || ((project left ty) =? 1)) || (((project left ty) =? 25) || ((project left ty) =? 26))) || ((project left ty) =? 30)) then (if (((project left ty) =? (project output ty)) && (repeatValid (project output n0) (project output n1) (project output n2) (project output n3) (project left n0) (project left n1) (project left n2) (project left n3))) then (match (TensorTypes::tensorTypeStorage (project left ty))
    | TensorTypes::TypeStorageDecision.unknownType => (ctor PointwiseDecision.unsupported)
    | TensorTypes::TypeStorageDecision.knownType spec => (ctor PointwiseDecision.copyBlocks (blockCount output 1) (project spec typeSize) (indexContext left 1) (indexContext output 1) (ctor CopyMode.repeated))) else (ctor PointwiseDecision.invalidShape)) else (ctor PointwiseDecision.unsupported)),
  executable: true
}

definition{
  name: concatExtent,
  parameters: [axis : Nat, dimension : Nat, left : Nat, right : Nat, output : Nat],
  result: Bool,
  body: (if (axis =? dimension) then (output =? (left + right)) else ((left =? right) && (output =? left))),
  executable: true
}

definition{
  name: planConcat,
  parameters: [
    axis : Nat,
    left : TensorDescriptor,
    right : (Option TensorDescriptor),
    output : TensorDescriptor
  ],
  result: PointwiseDecision,
  body: (match right
    | Option.none => (ctor PointwiseDecision.invalidShape)
    | Option.some value => (if (((axis <? 4) && (((project left ty) =? (project value ty)) && ((project left ty) =? (project output ty)))) && ((((concatExtent axis 0 (project left n0) (project value n0) (project output n0)) && (concatExtent axis 1 (project left n1) (project value n1) (project output n1))) && (concatExtent axis 2 (project left n2) (project value n2) (project output n2))) && (concatExtent axis 3 (project left n3) (project value n3) (project output n3)))) then (match (TensorTypes::tensorTypeStorage (project left ty))
      | TensorTypes::TypeStorageDecision.unknownType => (ctor PointwiseDecision.unsupported)
      | TensorTypes::TypeStorageDecision.knownType spec => (let block : Nat := (project spec blockSize); (if ((((prim remainder (project left n0) block 1 : Nat) =? 0) && ((prim remainder (project value n0) block 1 : Nat) =? 0)) && ((prim remainder (project output n0) block 1 : Nat) =? 0)) then (ctor PointwiseDecision.copyBlocks (blockCount output block) (project spec typeSize) (indexContext left block) (indexContext output block) (ctor CopyMode.concatenated axis (indexContext value block))) else (ctor PointwiseDecision.invalidShape)))) else (ctor PointwiseDecision.invalidShape))),
  executable: true
}

definition{
  name: numericPlan,
  parameters: [
    operation : Nat,
    variant : Nat,
    left : TensorDescriptor,
    right : (Option TensorDescriptor),
    output : TensorDescriptor
  ],
  result: PointwiseDecision,
  body: (match right
    | Option.none => (planPointwise operation variant (project left ty) 0 (project output ty) false (project left n0) (project left n1) (project left n2) (project left n3) 1 1 1 1 (project output n0) (project output n1) (project output n2) (project output n3))
    | Option.some value => (planPointwise operation variant (project left ty) (project value ty) (project output ty) true (project left n0) (project left n1) (project left n2) (project left n3) (project value n0) (project value n1) (project value n2) (project value n3) (project output n0) (project output n1) (project output n2) (project output n3))),
  executable: true
}

definition{
  name: blockOffset,
  parameters: [index : Nat, context : IndexContext],
  result: Nat,
  body: (let q0 : Nat := (prim quotient index (project context u0) 0 : Nat); (let q1 : Nat := (prim quotient q0 (project context n1) 0 : Nat); (TensorIndex::elementOffset (prim remainder index (project context u0) 0 : Nat) (prim remainder q0 (project context n1) 0 : Nat) (prim remainder q1 (project context n2) 0 : Nat) (prim quotient q1 (project context n2) 0 : Nat) (project context s0) (project context s1) (project context s2) (project context s3)))),
  executable: true
}

definition{
  name: repeatedBlockOffset,
  parameters: [index : Nat, source : IndexContext, destination : IndexContext],
  result: Nat,
  body: (let q0 : Nat := (prim quotient index (project destination u0) 0 : Nat); (let q1 : Nat := (prim quotient q0 (project destination n1) 0 : Nat); (rightOffset (prim remainder index (project destination u0) 0 : Nat) (prim remainder q0 (project destination n1) 0 : Nat) (prim remainder q1 (project destination n2) 0 : Nat) (prim quotient q1 (project destination n2) 0 : Nat) (project source u0) (project source n1) (project source n2) (project source n3) (project source s0) (project source s1) (project source s2) (project source s3)))),
  executable: true
}

definition{
  name: concatenatedOffsets,
  parameters: [
    index : Nat,
    left : IndexContext,
    right : IndexContext,
    destination : IndexContext,
    axis : Nat
  ],
  result: GraphResponse,
  body: (let q0 : Nat := (prim quotient index (project destination u0) 0 : Nat); (let q1 : Nat := (prim quotient q0 (project destination n1) 0 : Nat); (let i0 : Nat := (prim remainder index (project destination u0) 0 : Nat); (let i1 : Nat := (prim remainder q0 (project destination n1) 0 : Nat); (let i2 : Nat := (prim remainder q1 (project destination n2) 0 : Nat); (let i3 : Nat := (prim quotient q1 (project destination n2) 0 : Nat); (let edge : Nat := (if (axis =? 0) then (project left u0) else (if (axis =? 1) then (project left n1) else (if (axis =? 2) then (project left n2) else (project left n3)))); (let chosen : Bool := (edge <=? (if (axis =? 0) then i0 else (if (axis =? 1) then i1 else (if (axis =? 2) then i2 else i3)))); (let source : IndexContext := (if chosen then right else left); (ctor GraphResponse.offsets chosen (TensorIndex::elementOffset (if (chosen && (axis =? 0)) then (prim subtract i0 edge : Nat) else i0) (if (chosen && (axis =? 1)) then (prim subtract i1 edge : Nat) else i1) (if (chosen && (axis =? 2)) then (prim subtract i2 edge : Nat) else i2) (if (chosen && (axis =? 3)) then (prim subtract i3 edge : Nat) else i3) (project source s0) (project source s1) (project source s2) (project source s3)) (blockOffset index destination))))))))))),
  executable: true
}

definition{
  name: copyOffsets,
  parameters: [index : Nat, source : IndexContext, destination : IndexContext, mode : CopyMode],
  result: GraphResponse,
  body: (match mode
    | CopyMode.linear => (ctor GraphResponse.offsets false (blockOffset index source) (blockOffset index destination))
    | CopyMode.repeated => (ctor GraphResponse.offsets false (repeatedBlockOffset index source destination) (blockOffset index destination))
    | CopyMode.concatenated axis right => (concatenatedOffsets index source right destination axis)),
  executable: true
}

definition{
  name: evaluate,
  parameters: [request : GraphRequest],
  result: GraphResponse,
  body: (match request
    | GraphRequest.plan operation variant left right output => (ctor GraphResponse.planned (if (copyOperation operation) then (planStorage operation left (match right
      | Option.none => false
      | Option.some value => true) output) else (if (operation =? 19) then (planRepeat left output) else (if (operation =? 21) then (planConcat variant left right output) else (numericPlan operation variant left right output)))))
    | GraphRequest.offsets index source destination mode => (copyOffsets index source destination mode)),
  executable: true,
  production: {effects: [overflow], targets: [rust-core]}
}

theorem{
  name: paddedSourceOffset,
  parameters: [],
  statement: ((blockOffset 5 (record IndexContext {u0 := 2, n1 := 3, n2 := 1, n3 := 1, s0 := 4, s1 := 12, s2 := 36, s3 := 36})) = 28),
  proof: reflexivity{}
}

theorem{
  name: reshapedDestinationOffset,
  parameters: [],
  statement: ((blockOffset 5 (record IndexContext {u0 := 3, n1 := 2, n2 := 1, n3 := 1, s0 := 8, s1 := 32, s2 := 64, s3 := 64})) = 48),
  proof: reflexivity{}
}

theorem{
  name: repeatedFourAxisOffset,
  parameters: [],
  statement: ((repeatedBlockOffset 251 (record IndexContext {u0 := 2, n1 := 3, n2 := 2, n3 := 2, s0 := 4, s1 := 12, s2 := 40, s3 := 96}) (record IndexContext {u0 := 4, n1 := 9, n2 := 4, n3 := 4, s0 := 4, s1 := 16, s2 := 144, s3 := 576})) = 124),
  proof: reflexivity{}
}

theorem{
  name: concatInnerBoundary,
  parameters: [],
  statement: ((concatenatedOffsets 3 (record IndexContext {u0 := 2, n1 := 1, n2 := 1, n3 := 1, s0 := 4, s1 := 8, s2 := 8, s3 := 8}) (record IndexContext {u0 := 3, n1 := 1, n2 := 1, n3 := 1, s0 := 8, s1 := 24, s2 := 24, s3 := 24}) (record IndexContext {u0 := 5, n1 := 1, n2 := 1, n3 := 1, s0 := 4, s1 := 20, s2 := 20, s3 := 20}) 0) = (ctor GraphResponse.offsets true 8 12)),
  proof: reflexivity{}
}

theorem{
  name: concatOuterBoundary,
  parameters: [],
  statement: ((concatenatedOffsets 12 (record IndexContext {u0 := 2, n1 := 3, n2 := 2, n3 := 1, s0 := 4, s1 := 8, s2 := 24, s3 := 48}) (record IndexContext {u0 := 2, n1 := 3, n2 := 2, n3 := 2, s0 := 4, s1 := 8, s2 := 24, s3 := 48}) (record IndexContext {u0 := 2, n1 := 3, n2 := 2, n3 := 3, s0 := 4, s1 := 8, s2 := 24, s3 := 48}) 3) = (ctor GraphResponse.offsets true 0 48)),
  proof: reflexivity{}
}

theorem{
  name: concatEmptyLeft,
  parameters: [],
  statement: ((concatenatedOffsets 1 (record IndexContext {u0 := 0, n1 := 1, n2 := 1, n3 := 1, s0 := 4, s1 := 0, s2 := 0, s3 := 0}) (record IndexContext {u0 := 2, n1 := 1, n2 := 1, n3 := 1, s0 := 4, s1 := 8, s2 := 8, s3 := 8}) (record IndexContext {u0 := 2, n1 := 1, n2 := 1, n3 := 1, s0 := 4, s1 := 8, s2 := 8, s3 := 8}) 0) = (ctor GraphResponse.offsets true 4 4)),
  proof: reflexivity{}
}
\end{semanticmodule}
\end{lexlean}
